\section{Formalization, implementation and evidence}
\label{sec:evidence}

\subsection{Lean formalization}

The reference evaluator's correctness, the width bounds and a
carrier-cardinality bound with the exponential factors of \cref{thm:main}
have been formalized in Lean~4 with Mathlib \cite{Moura2021Lean4,mathlib2020};
the bit-time and bit-space clause of \cref{thm:main} is proved on paper in
\cref{sec:width} and has no formal counterpart. The formal statements follow
the sections above:
the partition (\lean{Meanders.FirstCrossing.firstHeightPartition}\texttt{firstHeightPartition}), the cut bijection
(\lean{Meanders.FirstCrossing.asymmetricCutEquiv}\texttt{asymmetricCutEquiv}), the forced prefix and FIFO joins
(\lean{Meanders.FirstCrossing.forcedPrefix_exact}\texttt{forcedPrefix\_exact}, \lean{Meanders.FirstCrossing.forcedFIFO_sound}\texttt{forcedFIFO\_sound}), the exterior cut
bijection (\lean{Meanders.FirstCrossing.exteriorArchCutEquiv}\texttt{exteriorArchCutEquiv}, \lean{Meanders.FirstCrossing.exteriorCount_cut}\texttt{exteriorCount\_cut}), the
key-level future equivalence (\lean{Meanders.FirstCrossing.stateFutureEquiv}\texttt{stateFutureEquiv}), the bijection between
accepted native words and connected sources (\lean{Meanders.FirstCrossing.acceptedSourceEquiv}\texttt{acceptedSourceEquiv}), the
paired layer invariant (\lean{Meanders.FirstCrossing.jointLayer_invariant}\texttt{jointLayer\_invariant}) and joint correctness
(\lean{Meanders.FirstCrossing.evaluateJoint_correct}\texttt{evaluateJoint\_correct}, \lean{Meanders.FirstCrossing.count_eq}\texttt{count\_eq}). The correctness statements
are proved against the independent public definitions of $\Closed$ and $\Open$
of the development (the combinatorial models of \cref{sec:prelim}), including
the conventions $\Closed(0)=0$, $\Open(0)=1$ and the parity dispatch
\eqref{eq:odd}. \Cref{app:lean} lists the supporting declarations for each
statement of this paper.

The formal resource statements are the width bounds
(\lean{Meanders.FirstCrossing.low_geometry_port_bound}\texttt{low\_geometry\_port\_bound}: at most $2(K-1)$ ports in \LOW;
\lean{Meanders.FirstCrossing.high_geometry_port_bound}\texttt{high\_geometry\_port\_bound}: at most $2(n-K)+2$ ports at every
stage of a \HIGH\ sector) and the carrier-cardinality bounds
\lean{Meanders.FirstCrossing.carrierCard_le}\texttt{carrierCard\_le},
\lean{Meanders.FirstCrossing.low_carrierCard_le}\texttt{low\_carrierCard\_le} and
\lean{Meanders.FirstCrossing.high_carrierCard_le}\texttt{high\_carrierCard\_le}, which bound the number of keys passing the
validity test of a stage by $(n+1)^{4}K\,4^{K-1}$ in \LOW\ and
$(n+1)^{4}(n-K+2)\,4^{n-K+1}$ in a \HIGH\ sector. These fix the exponential
factors $4^{K-1}$ and $4^{n-K}$ of \cref{thm:main}; their polynomial
prefactors are looser than those of \cref{prop:carrier}. Every successor
produced by the formal step passes that validity test and layers are merged
by key, so a layer never holds more keys than this bound, but the layer-level
inequality itself, the polynomial-factor bit-time and bit-space accounting of
\cref{thm:cost} and \cref{cor:balance}, and the sorting and arithmetic costs
are proved on paper only; no claim is made about compiled code, allocator
behaviour or resident memory. The development uses
only the standard axioms \texttt{propext}, \texttt{Classical.choice} and
\texttt{Quot.sound}, checked by an axiom audit.

\subsection{Implementation and certificates}

Two Rust evaluators implement the transfer. The \emph{reference} uses explicit
arrays, arbitrary-precision weights and runtime assertions of the invariants.
The \emph{production} evaluator packs a key into $128$ bits (the noncrossing
pairing as a Dyck word of at most $64$ ports and the four counters), performs
the surgery directly on the word through transition plans compiled once per
stage and counter block, keeps fixed-width weights that are widened in place
when a per-layer mass bound says a sum might overflow, and partitions each
layer into $256$ key-hashed buckets processed by a bounded number of workers
(at most one per bucket; small layers run on one worker); sectors whose port
bound exceeds $64$ fall back to the reference. Layer order and all integers
are designed to be independent of the worker count, and the shipped tests
check this at small orders for worker counts up to $257$; it is a tested
property, not a proved one. A Lean theorem
(\lean{Meanders.FirstCrossing.Packed.decode_wordSuccessors}\texttt{Packed.decode\_wordSuccessors}) proves that the packed codec refines the
reference transitions label by label; the production word surgery itself is
validated differentially, not verified.

Both evaluators can emit a sector certificate: the ordered sector metadata
and, per sector, the terminal ordinary total, the unbonused lower-return total
and the bonus, with either every sorted layer retained (full mode) or omitted
layers regenerated by the checker (compact mode). In both modes the checker
re-derives every legal successor and both channels layer by layer; digests are
integrity commitments only. The soundness theorem
(\lean{Meanders.Certify.FirstCrossingRun.check_joint_correct}\texttt{Meanders.Certify.FirstCrossingRun.check\_joint\_correct}) states that a
certificate accepted by this checker establishes both public counts. A
count-only manifest carries the claimed integers and no layer evidence;
checking it means re-running the reference evaluator. Four evidence classes
therefore arise for a reported value, from weakest to strongest:
\begin{enumerate}[nosep,label=(E\arabic*)]
\item \emph{production output only}: the integer returned by one run of a
  producer, without any independent check;
\item \emph{agreement with an independently known value}, here the OEIS
  terms of Jensen and Howroyd;
\item \emph{acceptance of a full or compact certificate by the compiled
  checker};
\item \emph{kernel replay of a concrete numerical instance}, carried out
  for $\Closed(2)=2$, $\Open(2)=1$ and $\Open(4)=3$.
\end{enumerate}
The shipped pipeline exercises (E3) through order $8$ and (E2) wherever a
pinned value exists. The production series of \cref{tab:campaign} consists
of count-only executions: its values within the published OEIS range are at
level (E2), and its values beyond that range are at level (E1).

\subsection{Validation}

\paragraph{Shipped checks.} The repository's tests and pipeline compare both
First-Crossing producers against the brute-force oracle at small orders and
at every threshold, verify full and compact certificates with the compiled
checker, kernel-replay the small instances above, and compare the producers
with the pinned OEIS values \cite{OEIS_A005315,OEIS_A005316}: $\Closed(n)$
for $1\le n\le28$ and $\Open(m)$ for $0\le m\le55$. ($\Closed(0)$ is
excluded from the comparison: this paper and the development use
$\Closed(0)=0$, whereas A005315 lists $1$ at index $0$; the open
convention $\Open(0)=1$ agrees with A005316.) The exterior identity of
\cref{lem:exterior} is checked on every matching and every cut through
$n=8$ ($33{,}097$ cases).

\paragraph{Development history (not shipped).} Before implementation the
mathematics was re-derived in an adversarial audit whose independent
constructor agreed with a literal enumeration of all ordered pairs through
$n=9$ on every candidate sector; during development the Rust evaluators also
agreed with an independent Python reference, with two earlier
connectivity-transfer evaluators through $n=14$, and with the literal
enumeration through $n=8$ at every threshold through $n=6$. These checks and
the earlier evaluators are author-reported history and are not part of the
repository; they cannot be rerun by a reader and are not counted as evidence
above.

\subsection{Measurements}

\Cref{tab:campaign} reports count-only wall time and peak resident memory of
the production evaluator in one fixed sweep: one repository revision, one
release binary, one machine (AWS EC2 \texttt{r7a.32xlarge}, $128$ vCPUs,
Linux x86-64, rustc~1.98), run on 2026-09-15 through
\texttt{scripts/evaluate.py} at every native order $0\le n\le30$, one fresh
process per order, count-only manifests and no in-process golden value. Each
run at $n\ge1$ returns $\Closed(n)$, $\Open(2n-1)$ and $\Open(2n)$; the run
at $n=0$ returns $\Closed(0)$ and $\Open(0)$ only. Every run records $96$
requested and effective workers in its metadata; small layers run on one
worker. Orders below $23$ each took under $7$ seconds and are omitted
from the table; the whole series took $36.6$ minutes of evaluator wall time.
The timing boundary is the child process: builds and any later verification
are excluded, while serialization and output of the count manifests are
included. These figures describe the measured range only and say nothing
about the asymptotic exponent, which is the subject of \cref{thm:cost}. The
run records (\texttt{meta.json} and the count manifests of every order) are
provided in the production-run archive accompanying the tagged release
described in \cref{app:repro}, rather than in the source tree.

\begin{table}[t]
\centering
\caption{Production evaluator on \texttt{r7a.32xlarge}, count only, one
fresh process per order and an effective worker limit of $96$ (small layers
run on one worker); every run also produced $\Open(2n-1)$ and $\Open(2n)$.
Memory is the reported peak resident set size in MiB.}
\label{tab:campaign}
\begin{tabular}{rrrrr}
\toprule
$n$ & workers (effective) & samples & wall time (s) & peak RSS (MiB) \\
\midrule
23 & 96 & 1 & 13.16 & 3578 \\
24 & 96 & 1 & 15.43 & 4388 \\
25 & 96 & 1 & 37.95 & 13788 \\
26 & 96 & 1 & 46.02 & 15187 \\
27 & 96 & 1 & 148.06 & 46097 \\
28 & 96 & 1 & 178.99 & 37599 \\
29 & 96 & 1 & 733.43 & 119100 \\
30 & 96 & 1 & 999.39 & 120072 \\
\bottomrule
\end{tabular}
\end{table}

The rows of \cref{tab:campaign} grow in a parity staircase. With
$K=\lfloor n/2\rfloor+1$ and $d=n-K$ we have $d(2m)=m-1$ and
$d(2m+1)=d(2m+2)=m$, so the \HIGH\ boundary width $2d+O(1)$ of
\cref{lem:highwidth} increases on every even-to-odd step and is unchanged on
every odd-to-even step, and the exponential factor $4^{d}$ of the carrier
bound of \cref{prop:carrier} quadruples or stays fixed accordingly (its
polynomial prefactor changes at every step). The measured peak
RSS of $36.7$\,GiB at $n=28$, $116.3$\,GiB at $n=29$ and $117.3$\,GiB at
$n=30$ is consistent with this staircase; the theorem bounds the number of
keys, not the resident memory of a process, and no more than consistency is
claimed. Measured wall-time ratios need not follow the exponential factor
either: from $n=28$ to $n=30$ the wall time grew by a factor of about $5.6$,
against $4$ for $4^{d}$ alone, and the polynomial factors of \cref{thm:cost}
together with implementation effects (bucket scheduling, memory traffic,
weight widening) account for such differences within the measured range.

The values $\Closed(n)$ for $1\le n\le28$ and $\Open(m)$ for $0\le m\le55$
are listed in OEIS \cite{OEIS_A005315,OEIS_A005316} (the first $24$ closed
and $44$ open terms are due to Jensen \cite{Jensen2000}, the remainder of
both tables to Howroyd). Every value of the series in that range agrees with
the listed term (level (E2) above; $\Closed(0)$ is excluded by convention as
explained under Shipped checks). One run at rank $r\ge1$ returns
$\Closed(r)$, its alias $\Open(2r-1)=\Closed(r)$, and $\Open(2r)$, so ranks
$27$ to $30$ returned
\begin{align*}
\Closed(27)&=\Open(53)=987{,}975{,}910{,}996{,}038{,}555{,}989{,}486,\\
\Open(54)&=2{,}594{,}805{,}632{,}585{,}289{,}523{,}975{,}474,\\
\Closed(28)&=\Open(55)=10{,}726{,}008{,}363{,}361{,}842{,}734{,}385{,}644,\\
\Open(56)&=28{,}235{,}899{,}288{,}344{,}793{,}178{,}333{,}732,\\
\Closed(29)&=\Open(57)=116{,}936{,}079{,}373{,}841{,}873{,}422{,}508{,}566,\\
\Open(58)&=308{,}496{,}356{,}015{,}441{,}090{,}010{,}351{,}094,\\
\Closed(30)&=\Open(59)=1{,}279{,}842{,}888{,}607{,}976{,}889{,}797{,}535{,}044,\\
\Open(60)&=3{,}383{,}262{,}281{,}572{,}596{,}355{,}695{,}881{,}934.
\end{align*}
Within the published OEIS range the sweep agrees term for term:
$\Closed(27)$, $\Closed(28)$ and $\Open(53)$ to $\Open(55)$ reproduce the
listed values. $\Closed(29)$, $\Closed(30)$ and $\Open(56)$ to $\Open(60)$
are not listed in either OEIS entry as of 2026-09-14. They are fixed-sweep
production outputs at level (E1): no independent numerical validation, no
second implementation result and no certificate replay is claimed for them.
The odd values $\Open(57)$ and $\Open(59)$ are aliases of $\Closed(29)$ and
$\Closed(30)$ written by the same run and provide no independent arithmetic
check. The values above make the computed Open range continuous through
$m=60$. Jensen's 2025 slides \cite{Jensen2025} analyse an open series of
about $60$ terms and show closed-data plots through $27$ terms, so comparable
values exist in unpublished form; whether any value above has been published
elsewhere has not been verified, and no priority is claimed.

At small orders the earlier (unshipped) connectivity transfers were faster
in wall time by orders of magnitude in development measurements. Such finite
measurements establish neither a practical crossover nor an asymptotic
wall-time claim; the contribution of this paper is the proved bound.
